\documentclass[11pt,a4paper]{article}
\usepackage[utf8]{inputenc}
\usepackage[T1]{fontenc}
\usepackage{lmodern}
\usepackage[french]{babel}
\usepackage[margin=1.9cm]{geometry}
\usepackage{microtype}
\usepackage{booktabs}
\usepackage{array}
\usepackage[table]{xcolor}
\usepackage{enumitem}
\usepackage[most]{tcolorbox}
\definecolor{navy}{HTML}{1F3864}
\definecolor{bluemid}{HTML}{2E5496}
\definecolor{lightblue}{HTML}{D6E0F0}
\definecolor{accent}{HTML}{C0491F}
\newtcolorbox{cle}[1][]{colback=lightblue!40,colframe=navy,boxrule=0.6pt,
  arc=2pt,left=7pt,right=7pt,top=5pt,bottom=5pt,before skip=6pt,after skip=6pt,#1}
\setlength{\parindent}{0pt}
\setlength{\parskip}{4pt}

\title{\vspace{-1.2cm}\color{navy}\bfseries Sept phrases à contredire\\[2pt]
  \large Votre expérience du terrain contre mes conclusions sur papier}
\author{Kélian Kaddouri}
\date{}

\begin{document}
\maketitle
\vspace{-1.1cm}

\begin{cle}
\textbf{Ce qu'on vous demande, et surtout ce qu'on ne vous demande pas.} \emph{Pas} de lire mon
mémoire. Je vous soumets sept phrases sur lesquelles ce que vous avez vu chez des clients vaut
mieux que n'importe quel calcul, et je vous demande d'essayer de les \textbf{contredire}. Un
désaccord argumenté m'est plus utile qu'un accord, et il sera cité comme tel. Comptez une
demi-heure.
\end{cle}

\textbf{De quoi il s'agit, en trois phrases.} Je cherche à chiffrer ce que coûte, en capital
réglementaire, le fait de mal appliquer le règlement européen sur la résilience informatique des
acteurs financiers. Mon modèle suppose qu'un problème sur un domaine \emph{en entraîne} un autre,
dans un certain ordre, plutôt que les problèmes ne surviennent chacun de leur côté. Cet ordre,
je n'ai pas pu le mesurer sur des bases de données : je l'ai déduit de rapports d'enquête publics,
et c'est exactement là que votre avis compte.

\textbf{Les cinq domaines dont je parle.} \textbf{Gouvernance} : qui décide, qui est responsable,
quels moyens la direction met sur la sécurité. \textbf{Incidents} : repérer un incident,
l'escalader, y répondre, le déclarer. \textbf{Tests} : tester avant de mettre en service, vérifier,
appliquer les correctifs. \textbf{Prestataires} : les fournisseurs et la dépendance à un service
extérieur. \textbf{Partage d'infos} : recevoir et faire circuler les alertes sur les menaces.

\begin{cle}[colback=accent!7,colframe=accent!70!black]
\textbf{À lire avant de répondre aux phrases 3 et 4.} Ces deux-là portent sur l'ordre dans lequel
les choses lâchent, et je les ai déduites de rapports d'enquête officiels. Or une enquête
officielle cherche toujours un responsable, et finit très souvent par conclure que la direction
n'encadrait pas assez. Il se peut donc que l'ordre que j'ai trouvé reflète la façon dont ces
rapports sont écrits, plutôt que ce qui se passe vraiment dans une entreprise.
\textbf{Votre réponse à ces deux phrases, fondée sur ce que vous avez vu et non sur des rapports,
est le seul moyen de faire la différence.} Si vous manquez de temps, répondez à celles-là.
\end{cle}

\section*{Les sept phrases}

\textbf{1.} Ces cinq domaines suffisent à décrire la solidité informatique d'un acteur financier.
Il n'en manque pas un, et ils ne se chevauchent pas au point de gêner.

\textbf{2.} Quand plusieurs domaines flanchent chez un même acteur, c'est souvent parce que l'un
a entraîné l'autre, et pas seulement parce qu'une cause extérieure commune les a touchés tous les
deux.

\textbf{3.} Une gouvernance défaillante dégrade les autres domaines, mais l'inverse ne se produit
pas : un incident, un test raté ou un prestataire qui tombe ne dégradent pas la gouvernance
elle-même. \emph{(La phrase que j'aimerais le plus vous voir contredire. Avez-vous déjà vu un
incident faire empirer la gouvernance ?)}

\textbf{4.} Les problèmes des autres domaines finissent par se voir dans la gestion des
incidents, mais une mauvaise gestion des incidents ne dégrade pas en retour les tests, les
prestataires ou la circulation de l'information.

\textbf{5.} Un prestataire partagé qui tombe abîme plusieurs domaines \emph{en même temps}, au
lieu de faire tomber les choses les unes après les autres. C'est différent des autres domaines, et
je le traite différemment.

\textbf{6.} Ces trois indicateurs sont ceux qu'une direction des risques suit vraiment, ou
pourrait produire sans grand chantier : le délai pour déclarer un incident, le temps qu'il faut
pour le circonscrire, et le degré de concentration des prestataires.

\textbf{7.} Il serait réaliste de demander aux acteurs financiers un registre d'incidents qui
note la date où l'incident \emph{a eu lieu} (et non celle où il a été déclaré), qui range chaque
incident par domaine plutôt que par type d'attaque, et qui compte un problème touchant vingt
acteurs comme \emph{un} événement à vingt victimes.

\section*{Votre réponse}

Pour chaque phrase : \textbf{D} d'accord \textbullet{} \textbf{N} d'accord mais avec une nuance
\textbullet{} \textbf{C} pas d'accord \textbullet{} \textbf{S} sans avis. Le commentaire est la
partie qui m'intéresse vraiment : un exemple précis vaut mieux qu'une appréciation générale.

\vspace{3mm}
\begin{center}\small
\renewcommand{\arraystretch}{2.7}
\setlength{\tabcolsep}{4pt}
\begin{tabular}{@{}c|c|c|c|c|p{9.4cm}@{}}
\toprule
 & \textbf{D} & \textbf{N} & \textbf{C} & \textbf{S} & \textbf{Commentaire, exemple vécu} \\
\midrule
1 & & & & & \\ \hline
2 & & & & & \\ \hline
\rowcolor{accent!8} 3 & & & & & \\ \hline
\rowcolor{accent!8} 4 & & & & & \\ \hline
5 & & & & & \\ \hline
6 & & & & & \\ \hline
7 & & & & & \\
\bottomrule
\end{tabular}
\end{center}

\vspace{3mm}
\textbf{La question la plus précieuse, pour finir.} Qu'est-ce qui manque ? Quelle phrase aurait dû
figurer dans cette liste, ou quel enchaînement voyez-vous régulièrement sur le terrain et dont je
ne parle nulle part ?

\vspace{16mm}\hrule\vspace{1mm}\hrule\vspace{1mm}\hrule

\vspace{4mm}
\begin{cle}
\textbf{Ce que je ferai de vos réponses.} Elles seront citées dans mon mémoire, et les
désaccords en premier. Si vous contredisez la phrase 3 ou la 4, je corrige le modèle : c'est
précisément pour cela que je vous écris. Merci de me dire votre rôle, votre secteur, depuis
combien de temps vous exercez, et si vous acceptez d'être cité par votre nom ou préférez rester
anonyme.
\end{cle}

\vspace{-1pt}
{\footnotesize\textcolor{navy}{\textbf{Merci du temps que vous y passerez.}} Kélian Kaddouri
\textbullet{} kkaddouri@nexialog.com}

\end{document}
